\section{The geometric counterexample}\label{sec:geometric-construction}

Let $k$ be a field of characteristic zero.  Consider the finite flat
degree-three marking morphism
\[
 \pi:\P^1\times\operatorname{Sym}^2(\P^1)
 \longrightarrow\operatorname{Sym}^3(\P^1),\qquad
 (x,Q)\longmapsto x+Q.
\]
Its ramification divisor is
\[
 D=\{(x,Q):Q\text{ contains }x\}.
\]
Choose $0,\infty\in\P^1$ and let $H\subset\operatorname{Sym}^3(\P^1)$ be
the hyperplane whose equation on the affine line
$\P^1\setminus\{\infty\}$ is
\begin{equation}\label{eq:zero-sum-hyperplane}
 x+y+z=0.
\end{equation}
On the small diagonal this becomes $3x=0$; after projective closure the
remaining intersection multiplicity two lies at $\infty$.  Thus $H$ is
tangent but not osculating to the twisted cubic of triple points.

Set
\[
 E=\pi^{-1}(H),\qquad
 U=\bigl(\P^1\times\operatorname{Sym}^2(\P^1)\bigr)\setminus(D\cup E),
 \qquad Y=\operatorname{Sym}^3(\P^1)\setminus H.
\]
Since $H$ is a hyperplane in $\P^3\simeq\operatorname{Sym}^3(\P^1)$,
$Y\simeq\A^3$.

\subsection{The auxiliary conic}

For fixed $x$, let $D_x$ and $E_x$ be the lines cut out by $D$ and $E$ in
$\{x\}\times\operatorname{Sym}^2(\P^1)$.  If
$\Gamma=\{2t:t\in\P^1\}$ is the Veronese conic, then $D_x$ is tangent to
$\Gamma$ at $2x$, while the $E_x$ form the pencil through
$p_\infty=2\infty$.  Define
\[
 \tau(x)=-x/2,\qquad \alpha=\tau^{-1},\qquad
 h(x)=D_x\cap E_x=x+\alpha(x),\quad \alpha(x)=-2x,
\]
where the last plus sign denotes addition of effective divisors.  The image
$C=h(\P^1)$ is a smooth conic, and the boundary lines become its chords
\begin{equation}\label{eq:auxiliary-chords}
 E_x=\overline{h(x)p_\infty},\qquad
 D_x=\overline{h(x)h(\tau(x))},
\end{equation}
with coincident endpoints interpreted as tangency
\cite{rezchikov-geometric}.

\begin{proposition}[affine source]\label{prop:geometric-affine-source}
The residual-intersection construction gives $U\simeq\A^3$.
\end{proposition}

\begin{proof}
For $(x,Q)\in U$, the line $\overline{h(x)Q}$ has a second intersection
\[
 \rho(x,Q)\in C.
\]
This varies regularly, with tangency counted twice.  By
\eqref{eq:auxiliary-chords}, the value $p_\infty$ occurs exactly on the
removed divisor $E$, so $\rho$ takes values in
$C\setminus\{p_\infty\}\simeq\A^1$.

For a residual point $s$, the removed divisor $D$ excludes exactly the one
marked point satisfying $s=h(\tau(x))$.  Hence
\[
 B=\{(x,s)\in\P^1\times(C\setminus\{p_\infty\}):
          s\ne h(\tau(x))\}
\longrightarrow C\setminus\{p_\infty\}
\]
is the complement of a section in a $\P^1$-bundle.  It is a torsor under a
line bundle; both the line bundle and torsor are trivial on $\A^1$, so
$B\simeq\A^2$.

For fixed $(x,s)\in B$, the point $Q$ varies on
$\overline{h(x)s}\setminus\{h(x)\}\simeq\A^1$.  This is again the
complement of a section in a $\P^1$-bundle.  Its associated line bundle and
torsor are trivial on $B\simeq\A^2$.  Therefore
\[
 U\longrightarrow B\longrightarrow\A^1,
 \qquad U\simeq\A^3.\qedhere
\]
\end{proof}

\begin{theorem}[geometric counterexample]\label{thm:geometric-counterexample}
After choosing affine-space identifications, the restriction of $\pi$ gives
an etale, generically three-to-one morphism
\[
 \A^3\simeq U\longrightarrow Y\simeq\A^3.
\]
Its Jacobian determinant in these coordinates is a nonzero constant.  In
particular it is not an isomorphism and, after base change to an algebraic
closure, is not injective.
\end{theorem}

\begin{proof}
The map $\pi$ is the finite flat cover that marks one point of an effective
divisor of degree three.  Its ramification divisor is $D$, so its restriction
to $U$ is etale.  A divisor with three distinct points has three markings;
hence the restriction has generic degree three.  Proposition
\ref{prop:geometric-affine-source} and the identification $Y\simeq\A^3$
give the displayed affine morphism.  Its Jacobian determinant is a unit in a
polynomial ring and is therefore a nonzero constant.
\end{proof}

This proves the existence of the counterexample without choosing polynomial
coordinates.  The next two sections give a compact explicit formula and then
certify, including at the projective root at infinity, that it is a coordinate
realization of this geometric morphism.
